\documentclass{situationist-preprint}

\title{Transparency and Its Limits}
\msid{EPI-15-a}
\msfamily{EPI-15}
\msdiscipline{Epistemology}
\mstype{Principal manuscript}
\mskeywords{self-knowledge; transparency; belief; answerability to reasons; recalcitrant attitudes; introspection; safety}

\begin{document}
\maketitle

\begin{abstract}
How does one know what one believes? Evans observed that, asked whether I believe there will be a third world war, I attend to the same considerations as I would if asked whether there will be a third world war. We defend the transparency account: one knows whether one believes $p$ by considering whether $p$. The puzzle is how attention to the world can yield knowledge of the mind. We solve it with the notion of answerability. An attitude is answerable to reasons when the subject's verdict on the question it settles fixes, and is not merely evidence of, whether she has it. For answerable belief, the self-ascription that follows the verdict is based on a state that partly constitutes the belief ascribed; we show that it is justified, safe and true, and so knowledge. The same argument extends to intention and to other attitudes that answer to a deliberative question. It fails for attitudes that are not answerable to reasons: the recalcitrant fear, the compulsive desire, the belief held against one's own verdict. For these the verdict and the attitude come apart, and the transparent self-ascription is unsafe. Self-knowledge of such attitudes must come by other routes, such as observation of one's conduct and introspection of how one feels.
\end{abstract}

\section{The problem}

The traditional account of self-knowledge is perceptual: one knows one's own mind by an inner sense that observes it, as one knows the world by outer senses. \citet{Shoemaker1994} argued against the model that, unlike a perceptual faculty, introspection could not be absent in a rational subject, and that there is no organ and no characteristic phenomenology of inner observation. An alternative begins from an observation of \citet[225]{Evans1982}: in making a self-ascription of belief, one's eyes are, so to speak, directed outward, upon the world. Asked whether I believe that $p$, I answer by putting into operation whatever procedure I have for answering the question whether $p$.

\citet{Moran2001} developed the observation into an account of first-person authority: the believer's relation to her own beliefs is that of an agent who makes up her mind, not of an observer who finds out. \citet{Byrne2005,Byrne2018} gave an epistemic version, on which one knows that one believes $p$ by following the rule: if $p$, believe that you believe that $p$. Following the rule requires recognizing that $p$, and so believing it, so a belief formed by following the rule is true. Critics have argued that the account covers only a small part of self-knowledge \citep{Gertler2011,Cassam2014}. Psychological work shows that people are often wrong about the causes of their attitudes \citep{NisbettWilson1977} and that their professed beliefs often diverge from those they act on \citep{Schwitzgebel2010,Gendler2008}.

We defend transparency as the method by which one knows whether one believes $p$, and we state precisely where the method fails. The puzzle of transparency is that the method seems to move from a proposition about the world to a proposition about oneself, which the first does not support \citep{Byrne2005}. We solve it by attending to what the self-ascription is based on, and the solution shows where the method stops working.

\section{Framework}

\begin{definition}[Transparent method]\label{d:method}
A subject follows the \emph{transparent method} for an attitude type $A$ with respect to a question $Q$ if she considers $Q$, reaches a verdict on it or fails to, and self-ascribes $A$ accordingly. For belief, $Q$ is whether $p$; she self-ascribes the belief that $p$ if her verdict is that $p$, and its absence otherwise.
\end{definition}

The method's question is the question that the attitude settles: for belief, whether $p$; for intention, whether to $\phi$; for fear, whether the object is dangerous; for admiration, whether it is admirable.

\begin{definition}[Answerability]\label{d:answer}
An attitude of a subject is \emph{answerable to reasons} if whether she has it is fixed by her verdict on the question it settles: she has it when her verdict is affirmative, lacks it when it is not, and it changes when her verdict changes.
\end{definition}

Answerability is a property of particular attitudes, not of attitude types. Most beliefs are answerable: a subject who concludes that the train has left believes it, and if she concludes on further evidence that it has not, she ceases to believe it. Some are not: the self-deceived parent who sincerely concludes, when he considers the evidence, that his son is guilty, but whose conduct, trust and expectations continue to be those of someone who believes him innocent. Fears and desires are often answerable, as when one ceases to fear the dog once one judges it harmless; but the phobic's fear of flying persists through her verdict that flying is safe, and is \emph{recalcitrant} \citep{Gendler2008}.

\begin{assumption}[Constitution]\label{a:const}
If a subject's belief about $p$ is answerable, her affirmative verdict on whether $p$ partly constitutes her believing $p$, and the absence of an affirmative verdict, on consideration, partly constitutes her not believing it.
\end{assumption}

Assumption~\ref{a:const} follows from Definition~\ref{d:answer} on a view on which a belief is a standing commitment to the truth of a proposition, of which the verdict is the occurrent manifestation. It does not say that every verdict creates a belief: a verdict that is not integrated with the subject's conduct, as in the self-deceived parent, is not partly constitutive of a belief, and the belief in that case is not answerable.

\begin{principle}[Knowledge]\label{p:know}
A subject knows $q$ if she believes $q$ on a basis on which she is justified, her belief is safe, and $q$ is true (EPI-01).
\end{principle}

\section{The argument}

The puzzle of transparency is solved by locating the basis of the self-ascription. The basis is not the proposition $p$, which is no evidence that anyone believes anything. It is the subject's own verdict that $p$, her judging that $p$, which is a mental state of hers and which she is in when she self-ascribes. Following the transparent method is moving from judging that $p$ to judging that one believes $p$, not from $p$ to the proposition that one believes $p$.

\begin{theorem}[Transparency]\label{t:trans}
If a subject's belief about $p$ is answerable, a self-ascription of belief or non-belief reached by the transparent method is knowledge.
\end{theorem}

\begin{proof}
Suppose her verdict is that $p$, and she self-ascribes the belief that $p$. By Assumption~\ref{a:const}, the verdict partly constitutes her believing $p$; so the self-ascription is true. It is safe: in close cases in which she self-ascribes the belief on the same basis, she has reached the verdict that $p$, and, the belief being answerable, the verdict partly constitutes the belief there too. It is justified: its basis is a mental state of hers which, given answerability, is the manifestation of the state ascribed, and on the internalist view of EPI-06 justification depends only on such states and her response to them; nothing she has to go on speaks against the self-ascription. By \ref{p:know}, she knows that she believes $p$. The case in which she reaches no affirmative verdict and self-ascribes non-belief is parallel, using the second clause of Assumption~\ref{a:const}.
\end{proof}

The theorem explains the features that made transparency attractive. The method is available to every reasoner, since it requires only the ability to consider whether $p$, and it needs no organ of inner sense. It explains first-person authority: the subject's self-ascriptions of answerable beliefs are not inferred from evidence about herself, and so are not open to challenge in the way third-person ascriptions are. And it explains Moore's paradox: one cannot sincerely assert \enquote{$p$, but I do not believe that $p$}, since the verdict that grounds the first conjunct grounds, by the method, the denial of the second.

\begin{corollary}[Extension]\label{c:ext}
For any attitude type that settles a question, a self-ascription of an answerable attitude of that type reached by the transparent method is knowledge.
\end{corollary}

\begin{proof}
The proof of Theorem~\ref{t:trans} uses only the answerability of the attitude and the parallel of Assumption~\ref{a:const} for its type, which follows in the same way from Definition~\ref{d:answer}: one knows whether one intends to $\phi$ by considering whether to $\phi$, and whether one fears the dog by considering whether it is dangerous, so long as one's intention or fear is fixed by one's verdict.
\end{proof}

\begin{theorem}[Limits]\label{t:limits}
If an attitude is not answerable to reasons, a self-ascription reached by the transparent method is not safe, and so is not knowledge, even when it is true.
\end{theorem}

\begin{proof}
If the attitude is not answerable, whether the subject has it is not fixed by her verdict. There are then cases, close to the actual one, in which her verdict on the question is the same and her attitude differs: her fear of flying persists, or would have persisted, through the verdict that flying is safe; the parent's belief in his son's innocence persists through the verdict of guilt. In those cases the method, applied on the same basis, yields a false self-ascription. So the self-ascription is unsafe, and by \ref{p:know} not knowledge. If in the actual case the attitude happens to agree with the verdict, the self-ascription is true, but true by luck relative to its basis.
\end{proof}

The phobic who asks herself whether she is afraid of flying, and answers by considering whether flying is dangerous, will conclude that she is not afraid; she is. The method yields the wrong answer whenever the attitude is recalcitrant, and an unreliable answer whenever recalcitrance is a close possibility. The failure is not a failure of attention or of reasoning, which may be impeccable; it is a mismatch between the method and the attitude, whose presence is not settled by the question the method asks.

\begin{corollary}[Other routes]\label{c:other}
Knowledge of attitudes that are not answerable to reasons, where it is possible, is acquired by methods other than transparency: by inference from one's conduct and reactions, by testimony and by introspection of their phenomenal character.
\end{corollary}

\begin{proof}
By Theorem~\ref{t:limits}, transparency does not yield it. The remaining sources are those by which one knows the attitudes of others, which apply to oneself as well \citep{Cassam2014}, and introspection of current experience. MND-20 holds that introspection is reliable about current phenomenal states and unreliable about their causes. The phobic can know by introspection that she feels fear when she boards, and by observing her conduct that she avoids planes; she cannot know by introspection why.
\end{proof}

EPI-15-b examines self-knowledge of recalcitrant attitudes in more detail.

\section{Objections}

\begin{objection}[Detection or formation]
Considering whether $p$ may form a new belief rather than detect an existing one. If so, the transparent method does not tell me what I believed, only what I now come to believe.
\end{objection}

\begin{reply}
For an answerable belief, considering whether $p$ on unchanged evidence reproduces the verdict that constitutes the standing belief; there is nothing further to detect. Where consideration changes the verdict, because the subject notices something she had not attended to, the belief has changed with it, and the self-ascription reports the present belief, which is what the question \enquote{do you believe $p$?} asks about. EPI-11 holds that belief is not under direct voluntary control but that attention is, and that belief follows attention; the transparent method is a deployment of attention, and its effect on belief is the effect that EPI-11 describes, not an act of will.
\end{reply}

\begin{objection}[Bad inference]
The transparent method is an inference from $p$ to \enquote{I believe that $p$}, which is invalid and not even inductively good. Byrne's rule is followed only by those who have the premise, but that does not make the inference good \citep{Gertler2011}.
\end{objection}

\begin{reply}
The method is not an inference from the proposition $p$; its basis is the subject's judging that $p$. A transition from a mental state to a belief that one is in a state it partly constitutes is not an inference from evidence, and it does not need the support of a premise connecting $p$ with one's beliefs. What makes it good is what Theorem~\ref{t:trans} shows: the transition is safe, and it is justified on the internal basis it has.
\end{reply}

\begin{objection}[Ubiquity of non-answerable belief]
Experimental work shows that professed beliefs often diverge from the beliefs manifested in action, as with people who sincerely judge that the races are equal in intelligence and act otherwise \citep{Schwitzgebel2010}. If such divergence is common, the transparent method fails for belief generally.
\end{objection}

\begin{reply}
The cases show that non-answerable belief-like states are common, which Theorem~\ref{t:limits} accommodates; they do not show that answerable beliefs are rare. For most of the propositions one has considered, whether the train has left, what one's address is, whether the meeting is on Tuesday, verdict and belief coincide, and the method yields knowledge. Where they diverge, the theorem predicts exactly what the experiments find: sincere self-ascriptions that are false. That the method fails where answerability fails is our thesis, not an objection to it.
\end{reply}

\begin{objection}[Rationality without self-knowledge]
A subject could reliably follow the method and still not know her own mind if she lacked the concept of belief. Transparency presupposes what it is supposed to explain.
\end{objection}

\begin{reply}
The method presupposes possession of the concept of belief, as every method of self-ascribing belief does; it does not presuppose knowledge of which beliefs one has, which is what it explains. A subject who has the concept and follows the method acquires that knowledge without any further self-observation.
\end{reply}

\section{Conclusion}

One knows whether one believes $p$ by considering whether $p$. The basis of the self-ascription is one's verdict, a mental state that, when the belief is answerable to reasons, partly constitutes the belief ascribed; the self-ascription is therefore justified, safe and true. The account extends to every attitude that settles a question and is fixed by one's verdict on it. Where an attitude is not answerable, verdict and attitude come apart, and the transparent self-ascription is unsafe. The recalcitrant fear, the compulsive desire and the belief held against one's own judgment must be known, if at all, as one knows the attitudes of others, or by introspecting how one feels.

\begin{thebibliography}{99}\small
\bibitem[Byrne(2005)]{Byrne2005} Byrne, A. (2005). Introspection. \emph{Philosophical Topics}, 33(1), 79--104.
\bibitem[Byrne(2018)]{Byrne2018} Byrne, A. (2018). \emph{Transparency and Self-Knowledge}. Oxford University Press.
\bibitem[Cassam(2014)]{Cassam2014} Cassam, Q. (2014). \emph{Self-Knowledge for Humans}. Oxford University Press.
\bibitem[Evans(1982)]{Evans1982} Evans, G. (1982). \emph{The Varieties of Reference}. Oxford University Press.
\bibitem[Gendler(2008)]{Gendler2008} Gendler, T.~S. (2008). Alief and belief. \emph{The Journal of Philosophy}, 105(10), 634--663.
\bibitem[Gertler(2011)]{Gertler2011} Gertler, B. (2011). Self-knowledge and the transparency of belief. In A.~Hatzimoysis (Ed.), \emph{Self-Knowledge}. Oxford University Press.
\bibitem[Moran(2001)]{Moran2001} Moran, R. (2001). \emph{Authority and Estrangement: An Essay on Self-Knowledge}. Princeton University Press.
\bibitem[Nisbett and Wilson(1977)]{NisbettWilson1977} Nisbett, R.~E. and Wilson, T.~D. (1977). Telling more than we can know: Verbal reports on mental processes. \emph{Psychological Review}, 84(3), 231--259.
\bibitem[Schwitzgebel(2010)]{Schwitzgebel2010} Schwitzgebel, E. (2010). Acting contrary to our professed beliefs, or the gulf between occurrent judgment and dispositional belief. \emph{Pacific Philosophical Quarterly}, 91(4), 531--553.
\bibitem[Shoemaker(1994)]{Shoemaker1994} Shoemaker, S. (1994). Self-knowledge and \enquote{inner sense}. \emph{Philosophy and Phenomenological Research}, 54(2), 249--314.
\end{thebibliography}

\end{document}
